% Hypothesis test options for Project 2, shared by Project2.tex and
% P2_1_activity.tex. Each document supplies its own section heading.

Pick the tests that match your variables and research questions: at least one categorical test and at least one numerical test.
Every test has an \textbf{independence} condition: the observations should come from a random sample (or a randomized experiment), and if sampling without replacement, the sample should be less than 10\% of the population.

\subsubsection*{Categorical Data Analysis (choose at least one)}
\begin{itemize}
    \item \textbf{Single proportion} (Chapter 6.1): one categorical variable with two outcomes. Compare the proportion in one outcome to a hypothesized value $p_0$.\\
    \emph{Example:} Is the proportion of Titanic passengers who survived different from 50\%?\\
    \emph{Conditions:} independence; success-failure using the null value: $np_0 \ge 10$ and $n(1-p_0) \ge 10$.
    \item \textbf{Difference of two proportions} (Chapter 6.2): a two-outcome variable compared between two groups. Test whether the proportion is the same in both groups.\\
    \emph{Example:} Is the survival rate different for female and male Titanic passengers?\\
    \emph{Conditions:} independence within each group and between the groups; success-failure in each group using the pooled proportion $\hat{p}_{pool}$ = (total successes in both groups)$/(n_1 + n_2)$: at least 10 expected successes and 10 expected failures in each group.
    \item \textbf{Chi-square goodness-of-fit} (Chapter 6.3): one categorical variable with three or more outcomes. Test whether the counts match a hypothesized distribution.\\
    \emph{Example:} Are births equally likely on all seven days of the week?\\
    \emph{Conditions:} independence; every expected count is at least 5; degrees of freedom $> 1$ (at least 3 categories).
    \item \textbf{Chi-square test of independence} (Chapter 6.4): two categorical variables, at least one with three or more outcomes. Test whether the two variables are associated.\\
    \emph{Example:} Is survival associated with passenger class (1st, 2nd, 3rd) on the Titanic?\\
    \emph{Conditions:} independence; every expected count is at least 5, where expected count = (row total)(column total)/(table total).
\end{itemize}

\subsubsection*{Numerical Data Analysis (choose at least one)}
For the $t$-tests below, the \textbf{normality} condition depends on the sample size: if $n < 30$, the data should be nearly normal with no clear outliers; if $n \ge 30$, there should be no extreme skew. Check it with a histogram or boxplot.
\begin{itemize}
    \item \textbf{One-sample mean} (Chapter 7.1): one numerical variable. Compare its mean to a hypothesized value $\mu_0$.\\
    \emph{Example:} Do college students sleep less than 8 hours per night on average?\\
    \emph{Conditions:} independence; normality of the data.
    \item \textbf{Paired means} (Chapter 7.2): two numerical measurements on the same individuals (e.g., before and after). Test whether the mean of the differences is zero.\\
    \emph{Example:} Do students score differently on the final exam than on the midterm, on average?\\
    \emph{Conditions:} the pairs are independent of each other; normality of the \emph{differences}.
    \item \textbf{Difference of two means} (Chapter 7.3): one numerical variable compared between two independent (not paired) groups. Test whether the two group means are equal.\\
    \emph{Example:} Is the average birth weight different for babies of mothers who smoke and mothers who don't?\\
    \emph{Conditions:} independence within each group and between the groups; normality in \emph{each} group.
    \item \textbf{ANOVA} (Chapter 7.5): one numerical variable compared across three or more groups. Test whether all the group means are equal.\\
    \emph{Example:} Does average gas mileage differ for cars with 4, 6, and 8 cylinders?\\
    \emph{Conditions:} independence within and between groups; the observations in each group are nearly normal (especially important for small groups); constant variance: the groups have roughly equal spread (compare the standard deviations or side-by-side boxplots).
\end{itemize}

\noindent\textbf{What about regression?} Questions about how one numerical variable changes with another numerical variable are answered with \textbf{linear regression}. Regression is \textbf{not required} for this project (it is not taught until MA214). You may still show a scatterplot of two numerical variables as part of your EDA, but your hypothesis test should be one of the tests above.\\
\emph{Example:} On average, how much does a car's gas mileage change for each additional 1{,}000 pounds of weight? Is this value significantly different from zero? (This is a regression question, not a difference of means question.)
